Este apêndice documenta o código-fonte desenvolvido em \textit{Python} para a engenharia de atributos e o treinamento das arquiteturas de \textit{Deep Learning} baseadas na biblioteca \textit{PyTorch}. Os scripts foram organizados conforme as etapas do \textit{pipeline} e estão disponíveis no diretório \texttt{Fonte/src/} do código-fonte anexado a este trabalho.

\section{Configuração do DataLoader e Normalização}
A função \texttt{prepare\_data\_dl} carrega o \textit{dataset}, aplica a normalização \texttt{MinMaxScaler} ajustada exclusivamente nos dados de treinamento, e a função auxiliar \texttt{create\_sequences} transforma a série temporal em janelas deslizantes de tamanho $L = 14$ (\textit{lookback}), gerando tensores tridimensionais $(B \times L \times F)$ compatíveis com as camadas recorrentes do PyTorch.

\lstinputlisting[language=Python, caption={Preparador de Dados (data\_loader\_dl.py)}, label={lst:dataloader}]{../Fonte/src/models/data_loader_dl.py}

\section{Rede Recorrente Bidirecional (LSTM-Bi)}
Implementação da arquitetura Bi-LSTM com processamento bidirecional e portas de esquecimento (\textit{Forget Gates}), utilizando Dropout ($p = 0{,}4$) e otimizador AdamW com \textit{Weight Decay} ($\lambda = 10^{-4}$) para regularização. O treinamento aplica \texttt{ReduceLROnPlateau} como \textit{learning rate scheduler} e \textit{gradient clipping} ($\|\nabla\| \leq 1{,}0$) para mitigar a explosão do gradiente característica de redes recorrentes profundas.

\lstinputlisting[language=Python, caption={Modelo Completo LSTM (lstm\_bidirecional.py)}, label={lst:lstm}]{../Fonte/src/models/lstm_bidirecional.py}

\section{Rede Gated Recurrent Unit Bidirecional (GRU-Bi)}
Arquitetura alternativa baseada em \textit{Gated Recurrent Units}, que simplifica a LSTM ao eliminar a célula de estado ($C_t$) separada. Reaproveita as funções utilitárias \texttt{reverse\_scaling} e \texttt{plot\_loss} do módulo Bi-LSTM, e a função genérica \texttt{evaluate\_and\_plot} do módulo XGBoost para padronizar a avaliação e a geração de gráficos. Aplica \texttt{ReduceLROnPlateau} como \textit{scheduler} e \textit{gradient clipping} idêntico ao da Bi-LSTM.

\lstinputlisting[language=Python, caption={Modelo Completo GRU (gru\_avancada.py)}, label={lst:gru}]{../Fonte/src/models/gru_avancada.py}

\section{Modelo XGBoost (\textit{Gradient Boosting})}
Implementação do XGBoost com previsão direta $h = 1$ (alvo = dia seguinte). Os hiperparâmetros foram selecionados por \textit{Grid Search} temporal (\texttt{TimeSeriesSplit}, 5 \textit{folds}, 288 combinações): $n\_estimators = 500$, $learning\_rate = 0{,}03$, $max\_depth = 4$, $subsample = 0{,}8$, $colsample\_bytree = 0{,}7$, $min\_child\_weight = 3$ e $\gamma = 0$, com \texttt{random\_state} fixado em 42 para reprodutibilidade. Os melhores parâmetros são exportados para \texttt{xgboost\_best\_params.json} e consumidos pela análise multi-horizonte. O módulo também define a função genérica \texttt{evaluate\_and\_plot}, reutilizada pelos modelos Bi-LSTM e Bi-GRU.

\lstinputlisting[language=Python, caption={Rotina do XGBoost (baseline\_xgboost.py)}, label={lst:xgboost}]{../Fonte/src/models/baseline_xgboost.py}

\section{Pipeline de Análise Exploratória (EDA)}
Script responsável pela execução automatizada da análise exploratória dos dados, incluindo a geração da matriz de correlação de Pearson, diagramas de dispersão vento--interrupções e \textit{boxplots} de sazonalidade mensal.

\lstinputlisting[language=Python, caption={Pipeline de Exploração (script\_exploration\_pipeline.py)}, label={lst:eda}]{../Fonte/src/models/script_exploration_pipeline.py}

\section{Gerador de Gráficos de Diagnóstico}
Script responsável pela geração dos gráficos comparativos de análise residual: distribuição KDE dos resíduos dos três modelos, recorte temporal de anomalia (El Niño 2023) e dispersão de heteroscedasticidade (erro absoluto vs.\ volume real). Consome os arquivos \texttt{predictions\_*.csv} produzidos por cada modelo durante a avaliação.

\lstinputlisting[language=Python, caption={Gerador de Gráficos (advanced\_plots.py)}, label={lst:plots}]{../Fonte/src/models/advanced_plots.py}

\section{Previsão Multi-Horizonte}
Script que implementa a avaliação dos três modelos nos horizontes de 1, 3, 7 e 14 dias (\texttt{Fonte/src/models/previsao\_multihorizonte.py}) com a estratégia de \textbf{previsão direta} (\textit{Direct Multi-Step}): para cada horizonte $h$, um modelo independente é treinado mapeando os atributos do dia $t$ diretamente para $\text{interrupcoes}[t+h]$, sem recursão e sem acumulação de erros. Todos os modelos (XGBoost, Bi-LSTM e Bi-GRU) recebem o mesmo conjunto de informações históricas disponível na data de origem $t$. A divisão treino/teste é feita pela data-alvo, com 365 pares de teste por horizonte (\texttt{01/06/2024}--\texttt{31/05/2025}). A listagem completa integra o pacote de código-fonte entregue.

\section{Geração dos Gráficos Multi-Horizonte}
Script complementar (\texttt{Fonte/src/models/plot\_multihorizonte.py}) que regenera os gráficos de degradação do erro (MAE e RMSE) por horizonte de previsão a partir do arquivo \texttt{previsao\_multihorizonte\_metricas.csv}, sem necessidade de re-executar os modelos. Produz três figuras: degradação do MAE com anotações manuais, degradação do RMSE e painel duplo MAE/RMSE. A listagem completa integra o pacote de código-fonte entregue.

\section{Scripts Complementares de Análise Exploratória}
Em complemento aos seis módulos acima, o diretório \texttt{Fonte/src/} contém quatro scripts dedicados à análise exploratória dos dados e à engenharia de atributos, cujo conteúdo é descrito no Capítulo~\ref{3_Metodologia} e no Capítulo~\ref{4_Resultados}:

\begin{itemize}
    \item \texttt{01\_eda\_sazonalidade.py}: decomposição STL aditiva e funções de autocorrelação (ACF/PACF) sobre a série diária de interrupções e precipitação.
    \item \texttt{02\_correlacoes\_nao\_lineares.py}: matrizes de correlação de Spearman e Kendall e correlogramas cruzados (\textit{cross-correlation}) entre variáveis climáticas e interrupções.
    \item \texttt{03\_feature\_engineering.py}: construção do \textit{dataset} final (\texttt{dataset\_engenharia\_features.csv}) a partir da base diária consolidada, aplicando defasagens, médias móveis exponenciais (EMA), desvio-padrão móvel e features cíclicas de calendário.
    \item \texttt{04\_eda\_basica.py}: geração da série temporal completa, distribuição estatística, evolução anual, comparação multi-escala de correlações e \textit{violin plot} térmico por tercis de severidade.
\end{itemize}

A listagem completa desses scripts foi omitida deste apêndice por concisão, mas integra o pacote de código-fonte entregue.
